\section{Discussion and outlook}\label{sec:discussion}
In this work, we studied systems that are driven out of equilibrium by a force which is non-conservative but can nevertheless be described by a Hamiltonian with a non-standard form. First, by assuming that the non-conservative forces do not modify the system's interaction with the heat bath, we have shown that such systems are generally not expected to relax to equilibrium despite their Hamiltonian structure.
We then presented a mapping between two distinct mechanisms that drive a system out of equilibrium: coupling different DoF to heat baths at different temperatures, and the presence of a specific type of Hamiltonian curl-forces. This allows for a mechanical interpretation of the origin of rotational motion often observed in systems coupled to multiple heat baths, such as the Brownian gyrator, as the result of deterministic non-conservative forces. The direction of gyration is qualitatively described by the curl of this force field, $\nabla\times\mathbf{F}$, which is obtained without solving the stochastic dynamics of the system. Such interpretation has the benefit of being intuitive and easier to grasp in comparison to alternative ways of qualitatively explaining such phenomena, such as through an analysis of the angular dependence of the diffusion tensor \cite{abdoli2025quadrupolar}. It is clear that $\nabla\times\mathbf{F}$ cannot fully describe the currents in the thermal system as it lacks information regarding the steady state distribution itself. However, we note that other quantities suggested for the characterization of the non-trivial dynamics arising in such systems, such as the curl of the probability current \cite{dotsenko2013two}, are not necessarily more informative, as they tend to be complicated and do not carry a clear physical significance. Hence, in some cases $\nabla \times \mathbf{F}$, which is easily computed for arbitrarily complex systems, may be sufficient to describe some properties of the system. We further note that if $\nabla \times \mathbf{F} = 0$ in the curl-force representation, then the system is in thermal equilibrium with the bath. In the multiple-bath description, this is reflected in each degree of freedom independently equilibrating with its own bath.

The exactly solvable case of the OU process allowed us to explicitly compute the energetics of the system with curl-forces. In the steady state, we found that positive work is always applied on the system, regardless of the direction of gyration. This work is fully dissipated into the single heat bath as heat, and is associated with a positive EPR. Notably, while the mechanism driving the system out of equilibrium in each system is different (heat flowing between two reservoirs, compared to work being dissipated in the curl-force system), the EPR\textcolor{blue}{---an information theoretic quantity---}of the two systems is equal. Nevertheless, the two systems are distinguishable from one another, because experimentally measurable quantities such as the dissipated heat are not the same in the two descriptions.
We note that in contrast to gyrators, where heat flows from the bath with the higher temperature to that with the lower one, in the curl-force system work is always applied on the system and dissipated into the single bath, a consequence of the second law of thermodynamics.

% A powerful implication of the presented transformation is the possibility of making non-trivial qualitative predictions on systems for which an exact solution is unknown, as discussed in the example of the gyrator on a ring. In principle, the same analysis of $\nabla\times \mathbf{F}$ can be applied to other, arbitrarily complicated force fields, $\mathbf F$, to gain insights on the currents expected in a NESS (if such exists).

More generally, it will be interesting to further investigate the role of $\nabla \times \mathbf{F}$ on the nonequilibrium dynamics and energetics of a driven system. We have shown that not only it can be used to predict the direction of circulating currents in a NESS in both cases we studied, but also that in the exactly solvable linear case, where $\nabla\times\mathbf{F}$ is a constant, quantities such as power, heat dissipation rate and EPR are a function of its magnitude. Whether and how this result generalizes to other systems remains an open question.

One especially promising application of our method may be in experimental realizations of systems with multiple heat baths. Some experiments have been performed in the past, for example using system analog to Brownian motion consisting of an electric circuit with resistors held at different temperatures \cite{ciliberto2013heat, chiang2017electrical}. There have also been experiments conducted directly on Brownian particles, where in this case the anisotropic noise is artificially implemented by an athermal random force \cite{argun2017experimental, berut2014energy, lin2022stochastic, berut2016stationary, martins2026role}. However, none of these experiments featured diffusion with ``real'' thermal anisotropic noise, because it is not clear what it means for a single Brownian particle to be coupled to two heat sources simultaneously. On the contrary, curl-forces may be readily generated in the lab with existing setups, for example using optical tweezers \cite{roichman2008influence,roichman2008optical, sun2009brownian, wu2009direct, liberal2013near, sukhov2017non}. This could pave the way for an easier approach to effectively simulating experiments with multiple temperatures by performing experiments on the curl-force description instead, before transforming results back to the multiple-temperatures description.

This work was limited to a small subset of curl-forces that are described by a Hamiltonian with anisotropic---yet quadratic and positive---kinetic energy \cite{berry2015hamiltonian}. Generalizing to other types of non-conservative forces or Hamiltonians \cite{yip2024general}, as well as non-mechanical thermodynamic forces, could possibly lead to mappings of other types of nonequilibrium systems beyond the gyrator and its immediate generalizations. In essence, applying a similar canonical transformation to arbitrary nonequilibrium stochastic systems could lead to a more general form of ``nonequilibrium Hamiltonians'', which could eventually help form a unified theory for nonequilibrium systems \cite{ding2022unified}. One promising direction in this respect is performing a more rigorous mathematical study of curl-force dynamics and the identification of new conserved quantities, for example as was recently done by Yavari and Goriely \cite{yavari2025darboux}. While the conserved quantity they identified (specifically, a Hamiltonian) is only defined implicitly in terms of an integral equation so it cannot be applied directly to any system with curl-forces, perhaps explicit solutions could be derived for specific systems. In combination with our method, this could be a step towards finding closed form NESS solution for certain systems with multiple heat baths beyond the exactly solvable OU process. Finally, it is interesting to note that the Brownian gyrator can also be mapped to an information engine \cite{leighton2024information}. It is worthwhile considering how the curl-force description of the system relates to the feedback-control one, and more generally how these two methods of driving a system out of equilibrium relate to one another.

Other extensions of this work are also possible. It is interesting to consider whether this framework can be extended to the overdamped regime. The Brownian gyrator was originally formulated in the this limit \cite{filliger2007brownian}, but within our framework we are unable to give a mechanical interpretation of the rotational motion in this limit, as our description relies on the definitions of momentum and kinetic energy. Another possible application is studying the curl-force representation of gyrators with potentials other than the harmonic or ring potentials studied here, such as those considered in \cite{chang2021autonomous}, which could possibly provide more general insights. Finally, it is intriguing to construct a kinetic theory for a closed system consisting of particles with a non-isotropic kinetic energy term, namely when the bath itself has a curl-force Hamiltonian. Would such a system relax to the corresponding equilibrium, or to a nonequilibrium steady state? 